\documentclass[11pt,a4paper]{article}
\usepackage[T1]{fontenc}
\usepackage[utf8]{inputenc}
\usepackage{mathpazo}
\usepackage[margin=25mm]{geometry}
\usepackage{amsmath,amssymb,amsthm,mathtools,booktabs,array,longtable}
\usepackage{microtype,xcolor,enumitem,listings,fancyhdr,xurl,hyperref}
\hypersetup{colorlinks=true,linkcolor=blue!45!black,urlcolor=blue!45!black,
 citecolor=blue!45!black,pdftitle={Thirteen Integer Variables through a Single Five-Radical Norm},
 pdfauthor={Research contribution for Vladimir Reshetnikov's ProveIt project}}
\setlength{\headheight}{14pt}\pagestyle{fancy}\fancyhf{}
\fancyhead[L]{\small Thirteen integer variables}
\fancyhead[R]{\small ProveIt research contribution}\fancyfoot[C]{\thepage}
\setlength{\parskip}{4pt}\setlength{\parindent}{0pt}
\setlist{nosep,leftmargin=1.7em}\emergencystretch=2em
\newtheorem{theorem}{Theorem}[section]
\newtheorem{lemma}[theorem]{Lemma}
\newtheorem{proposition}[theorem]{Proposition}
\theoremstyle{definition}\newtheorem{definition}[theorem]{Definition}
\theoremstyle{remark}\newtheorem{remark}[theorem]{Remark}
\newcommand{\Z}{\mathbb Z}\newcommand{\Q}{\mathbb Q}\newcommand{\N}{\mathbb N}
\newcommand{\Con}{\operatorname{Con}}\newcommand{\Refutation}{\operatorname{Ref}}
\newcommand{\F}{\mathcal F}\newcommand{\Cal}{\mathcal K}
\newcommand{\pos}{\operatorname{Pos}}\newcommand{\Pair}{\operatorname{Pair}}
\newcommand{\Dm}{D_{\rm m}}\newcommand{\Dt}{D_{\tau}}
\newcommand{\Df}{D_f}\newcommand{\Dy}{D_y}\newcommand{\Ag}{A_{\rm g}}
\newcommand{\Cn}{C_{\rm n}}\newcommand{\Rtop}{R_{\rm top}}
\numberwithin{equation}{section}
\lstset{basicstyle=\ttfamily\small,breaklines=true,columns=fullflexible,
 keepspaces=true,frame=single,rulecolor=\color{black!20}}
\title{\vspace{-10mm}\textbf{Thirteen Integer Variables\\
through a Single Five-Radical Norm}\\[3mm]
\large Explicit consistency candidates, ordinary proofs,\\
and a precise calibration boundary}
\author{Research contribution prepared for\\Vladimir Reshetnikov's ProveIt project}
\date{24 October 2026\\\small Repository snapshot:
 \texttt{9a33342a03d19cf5b061b2083c6b5d3f27f752ed}}
\begin{document}
\maketitle
\begin{abstract}
We give a completely specified polynomial in thirteen unrestricted integer
variables, of exact total degree \(1088\), evaluated by a division-free circuit
with \(195\) live binary arithmetic operations. A simpler companion has
degree \(1392\) and \(186\) operations. Both use the six fixed
coefficients already present in ProveIt's source-34/source-71 construction.
There are no free parameters, variable exponents, hidden existential
coordinates, or new coefficient-majorant payloads.

The new mechanism places a guard inside a radical norm. Four radicands are
positive before any Diophantine condition is known. A zero of the
five-radical norm therefore forces its remaining, guard radicand to be
nonnegative. That fact supplies strict positivity inequalities; those
inequalities then make the five radical magnitudes strictly
superincreasing. Rationality and integrality follow, restoring four
eliminated square witnesses and the main Pell congruence.
An integer translation lowers the degree from \(1408\) to \(1392\).
A second thirteen-variable construction uses one carrier for two coprime
congruences to reduce the exact degree further to \(1088\), with \(195\)
binary operations. We retain the degree-\(1392\) formula as the simpler
base presentation and prove the additional CRT variant separately.

The algebraic reduction is proved here. Its interpretation as a
consistency equation composes explicitly identified, inherited ordinary
native-compiler and K215 calibration results. At that ordinary theorem
level, ZFC proves that the polynomial has no integer zero if and only if
\(\Con(\mathrm{ZFC})\). This contribution does not supply a new
kernel-checked calibration theorem or the repository's still-open
encoded object-ZFC derivation. It proves neither consistency nor global
minimality. The archive also supplies lower-degree companions,
coefficient reconstruction, exact norm and degree audits, and editable sources.
\end{abstract}
\clearpage
\begingroup\small\setlength{\parskip}{0pt}
\tableofcontents
\endgroup
\clearpage

\section{Results and proof status}

The request for a simplest consistency equation has several competing
objectives: unrestricted integer unknowns, total degree, circuit size,
coefficient size, and completeness of the proof certificate. We report
explicit tradeoffs. Every variable in Table~\ref{tab:new} ranges over \(\Z\).
\begin{table}[ht]\centering
\begin{tabular}{lrrr}\toprule
Construction & Unknowns & Exact degree & Binary operations\\\midrule
Five-root norm with coprime carrier &13&1088&195\\
Five-root norm, translated &13&1392&186\\
Five-root norm, unshifted &13&1408&185\\
Four-root norm, centered &14&696&150\\
Four-root norm, unshifted &14&792&137\\
Three-root norm, centered &15&344&131\\
Quotient fusion &20&24&120\\
Quotient fusion and transport &19&28&116\\
Quotient fusion and transport &18&32&112\\
Quotient fusion and transport &17&88&110\\\bottomrule
\end{tabular}
\caption{New literal equations. Operations are live binary \(+,-,\times\)
gates with shared identical subexpressions; fixed numeral leaves are not
gates. This does not measure bit complexity or expanded monomial count.}
\label{tab:new}\end{table}

The inspected repository's earlier thirteen-integer construction has
degree \(10\,461\,744\). On that specific variable/degree comparison,
\(1392\) reduces the degree by a factor exceeding \(7500\), and
\(1088\) by a factor exceeding \(9600\).
This is not a variable-count record: incoming material already gives
a ten-integer construction of degree \(1\,245\,324\).
Nine \emph{positive} variables constitute a different domain
comparison~\cite{repo,frontier,incoming}. Sun's recent ten-unknown
theorem is also relevant~\cite{sun2026}.

The formal product of conjugate radicals is a classical
relation-combining technique. The new feature is the absorption of an
indeterminate-sign guard without a circular proof or large separating
weights. The determinant guard supplies its own quantitative bound.

\subsection{Three distinct proof levels}
\begin{enumerate}
\item \textbf{New algebraic theorem.} The thirteen-variable polynomial is
existentially equivalent to the explicit normalized system in
Section~\ref{sec:kernel}. Its denominator argument, square extraction,
and exact degree are proved here.
\item \textbf{Inherited ordinary calibration.} The source-34 compiler,
signed/residue-oriented native theorem, canonical Pell construction,
and ordinary K215 theorem identify the relevant witnesses with finite
syntactic ZFC refutations. We use these named ordinary results as
imported mathematical lemmas.
\item \textbf{Formal certificate status.} Finite checks establish
identities, coefficients, circuits, and degree certificates. They do not
prove an unbounded compiler theorem. The repository separately records
that the complete native calibration and actual encoded
\texttt{BProv ZFCax\_s [] (fIff \ldots)} endpoint remain open
formalization obligations~\cite{boundary,calibration}.
\end{enumerate}

Thus a claim of ZFC provability below means ordinary mathematical
provability by composition with the specified ordinary proofs.
It does not mean that a completed Lean/Rocq or encoded object-language
derivation is delivered.

\section{The fixed polynomial}\label{sec:equation}

The thirteen unknowns, in literal circuit order, are
\[
\boxed{j,\ C,\ z,\ h,\ r,\ v,\ w,\ T,\ c,\ n,\ t,\ p_2,\ p_3\in\Z.}
\]
All other letters are polynomial abbreviations; \(C\) and \(c\) are
different unknowns. Use the complete
\texttt{mixed-parity-unrounded-1064921} tuple
\[
\beta=2^{1064921}-1,\qquad K_0,\quad M_C,\quad M_F,
\]
with the \emph{already shifted} \(M_F=M_{F,\rm native}+\beta\), and
\[
A_{\rm cone}=2^{1113},\qquad B_{\rm cone}=2^{666104}.
\]
Appendix~\ref{app:coeff} specifies the finite coefficient recipe.
The archive contains all six exact numeral payloads, the complete tile
and machine data, and an independent reconstruction. No coefficient
is selected by a halting search or a consistency assumption.

Define the following registers:
\begin{align}
J&=6j+1,&q&=\beta J+1,&W&=12w+4,\nonumber\\
X&=qW,&Y&=(2v+1)q^3,&E&=XY,\nonumber\\
a&=E+Y,&H&=4a+3,&\Delta&=a^2+H=(a+2)^2-1.
\label{eq:front1}
\end{align}
The transport equation is built into
\begin{align}
F&=(K_0+W)C-z(q-1),&\alpha&=q-F-C,\nonumber\\
R&=\bigl(q(q-F)-C\bigr)(q^2-1)+(M_C+qM_F)J,\nonumber\\
k&=R+1+hE,&\eta&=c-Yk,&\zeta&=k-\eta.
\label{eq:front2}
\end{align}
The abbreviation \(\alpha\) is only used in the proof. Set
\begin{align}
I&=(2r+1)c^2,&\Df&=1+\Delta I^2,&Q&=\Delta^2I^2,\nonumber\\
\delta&=Q-1,&V&=\Df(cT-R)-c,&C_y&=QV^2-1,\nonumber\\
\Dt&=1+XY^2(XY^2+1)k^2,&
\Dm&=1+\Delta c^2,&b_0&=X+ac.
\label{eq:front3}
\end{align}
Finally define
\begin{align}
U&=A_{\rm cone}F,&Z_c&=B_{\rm cone}C,\nonumber\\
G&=UZ_c-(q-U-Z_c)^2,&P_\eta&=\eta\zeta,\nonumber\\
\Ag&=GP_\eta-(t^2-G-P_\eta+1)^2
-1-p_2^2-p_3(p_3+1).
\label{eq:guard}
\end{align}

Let the universal five-root polynomial be
\begin{equation}
\F_{32}(x;A_0,\ldots,A_4)=
\prod_{\varepsilon\in\{-1,1\}^5}
\left(x-\sum_{i=0}^4\varepsilon_i\sqrt{A_i}\right).
\label{eq:F32}
\end{equation}
Changing any formal radical's sign permutes the factors, so the product
belongs to \(\Z[x,A_0,\ldots,A_4]\). Each radicand has degree at most
\(16\). Our equation is the cleared polynomial
\begin{equation}
\boxed{\;
P_{13}=\delta^{16}\F_{32}\!\left(
H(n+V)+b_0;\
\Dm,H^2\Df,H^2\Dt,H^2\Ag,\frac{H^2C_y}{\delta}
\right)=0.\;}
\label{eq:main}
\end{equation}
The denominator is canceled \emph{symbolically}. No division or
square-root operation occurs in the actual equation.
Section~\ref{sec:circuit} gives a recurrence using only integer
polynomial arithmetic, and
\texttt{equations/norm13\_degree1392.json} supplies every circuit gate.

Replacing \(n+V\) by \(n\) gives the degree-\(1408\), \(185\)-operation
companion. The two zero sets are related by the invertible integral
translation \(n_{\rm old}=n_{\rm new}+V\); \(V\) does not involve \(n\).

\section{A division-free norm circuit}\label{sec:circuit}

The three-root norm is
\[
\F_8(x;A,B,C)=
\bigl((x^2-A-B-C)^2-4(AB+AC+BC)\bigr)^2-64ABCx^2.
\]
Its cleared version is given by
\begin{align}
J_N&=\delta(x^2-A-B)-\Cn,\nonumber\\
T_N&=\delta^2AB,&R_N&=\delta\Cn,\nonumber\\
K_N&=J_N^2-4\bigl(T_N+R_N(A+B)\bigr),\nonumber\\
N_3(x;A,B,\Cn,\delta)&=K_N^2-64T_NR_Nx^2.
\label{eq:cleared3}
\end{align}
This is \(\delta^4\F_8(x;A,B,\Cn/\delta)\) after clearing.

To obtain the five-root norm, calculate \(N_3(x-u-v)\) in the free
rank-four algebra with basis \(1,u,v,uv\), \(u^2=D\), \(v^2=B_g\).
Squaring \((a_0,a_1,a_2,a_3)\) gives
\begin{align}
b_0&=a_0^2+Da_1^2+B_ga_2^2+DB_ga_3^2,\nonumber\\
b_1&=2(a_0a_1+B_ga_2a_3),&
b_2&=2(a_0a_2+Da_1a_3),\nonumber\\
b_3&=2(a_0a_3+a_1a_2).
\label{eq:biquad}
\end{align}
The square of \(x-u-v\) is
\((x^2+D+B_g,-2x,-2x,2)\).
If the resulting \(N_3\) has vector \((f_0,f_1,f_2,f_3)\), set
\begin{align}
g_0&=f_0^2+B_gf_2^2-D(f_1^2+B_gf_3^2),\nonumber\\
g_1&=2(f_0f_2-Df_1f_3),&
N_5&=g_0^2-B_gg_1^2.
\label{eq:norm5circuit}
\end{align}
Exactly,
\[
N_5=\delta^{16}\F_{32}(x;A,B,D,B_g,\Cn/\delta).
\]
For \eqref{eq:main} use \(A=H^2\Df\), \(B=H^2\Dt\),
\(D=\Dm\), \(B_g=H^2\Ag\), \(\Cn=H^2C_y\).
The identity holds with zero, negative, or dependent radicands;
the algebra need not be a field of degree \(32\).
The independent audit constructs the conjugate product in reverse
elimination order and compares all \(20\,349\) nonzero coefficients,
including the cleared dictionary.

\section{The guard establishes its own hypotheses}\label{sec:bootstrap}

\begin{lemma}[Unconditional realness]\label{lem:realness}
For every integer input, \(\delta\ne0\), \(H\ne0\), and
\[
\Dm>0,\quad\Df>0,\quad\Dt>0,\quad \Dy:=C_y/\delta\ge1.
\]
\end{lemma}
\begin{proof}
The fixed \(\beta\) is odd, \(J=6j+1\ne0\), and
\(|q|\ge\beta-1\ge16\). Thus \(q\) is even.
Since \(W=4(3w+1)\ne0\),
\(|X|\ge64\), \(|Y|\ge4096\), and
\(|a|\ge(|X|-1)|Y|>18\).
Consequently \(\Delta=(a+2)^2-1>1\).
Also \(a\) is even, so \(\Delta\) and \(H\) are odd.
The first two radicands are positive, and
\(\Dt\ge1\) because \(u(u+1)\ge0\) for the integer \(u=XY^2\).

If \(c=0\), then \(Q=0,\delta=-1,C_y=-1,\Dy=1\).
If \(c\ne0\), then \(Q>1\) and
\(\Df=1+\Delta(2r+1)^2c^4>|c|+1\).
For \(b=cT-R\), \(b=0\) gives \(|V|=|c|\ge1\);
otherwise \(|V|\ge\Df-|c|>1\).
Hence
\[
\Dy=V^2+\frac{V^2-1}{Q-1}\ge1.
\]
Both cases exclude \(\delta=0\).
\end{proof}

\begin{lemma}[Joint positivity]\label{lem:pair}
Put \(\Pair(x,y,t)=xy-(t^2-x-y+1)^2\).
For integers \(x,y\), some \(t\) has \(\Pair(x,y,t)>0\)
if and only if \(x,y>0\).
Every positive integer is a value of
\[
\pos(p_1,p_2,p_3)=1+p_1^2+p_2^2+p_3(p_3+1).
\]
\end{lemma}
\begin{proof}
Positivity of \(\Pair\) gives \(xy>0\).
If \(x=-a,y=-b\) with \(a,b>0\), then
\((t^2+a+b+1)^2>ab\), a contradiction.
Conversely let \(N=x+y-1\ge1\); \(xy\ge N\).
For \(m=\lfloor\sqrt N\rfloor\), choose \(t=m\) if \(N-m^2\le m\),
and \(t=m+1\) otherwise. Then \(|t^2-N|\le m\).
The error is zero when \(N\) is square, and otherwise its square
is less than \(N\le xy\).

For \(s>0\), \(4s-3\) is \(1\) or \(5\bmod8\), so Legendre's
three-square theorem applies. Modulo \(4\), exactly one square is
odd. Writing the squares as
\((2p_1)^2,(2p_2)^2,(2p_3+1)^2\) gives the formula for \(s\).
This is the inherited joint-guard lemma~\cite{joint}.
\end{proof}

\begin{proposition}[Realness before squarehood]\label{prop:guard}
At every integer zero of \(P_{13}\),
\[
\Ag\ge0,\qquad G>0,\qquad\eta\zeta>0,\qquad c\ne0,\quad Q>1.
\]
\end{proposition}
\begin{proof}
Lemma~\ref{lem:realness} permits interpretation of the cleared norm
over \(\mathbb C\). A linear factor vanishes. Four radicals and the
center are real. If \(\Ag<0\), the fifth radical
\(H\sqrt{\Ag}\) is the only nonzero imaginary term, which is
impossible. Thus \(\Ag\ge0\).
Now
\[
\Pair(G,\eta\zeta,t)=\Ag+1+p_2^2+p_3(p_3+1)>0.
\]
Lemma~\ref{lem:pair} gives both strict guards, without first assuming
that \(\Ag\) is square. If \(c=0\), then
\(\eta\zeta=-Y(Y+1)k^2\le0\), a contradiction.
\end{proof}

This order avoids a circular argument: realness establishes the guards;
the guards will establish separation; separation will establish squarehood.

\section{Automatic separation of five radicals}\label{sec:separation}

All bounds here precede Pell-index or tableau decoding. Put
\[
x=|X|\ge64,\quad m=|Y|\ge4096,\quad K=|k|,\quad C_0=|c|.
\]
Since \(\eta,\zeta\) are nonzero integers of the same sign,
\(K\ge2\), \(0<\theta=\eta/k<1\), and \(c=k(Y+\theta)\). Thus
\begin{equation}
K(m-1)<C_0<K(m+1).
\label{eq:cstrip}
\end{equation}
Write \(f_+=\sqrt{\Df}\), \(\tau_+=\sqrt{\Dt}\),
\(y_+=\sqrt{\Dy}\), \(d_+=\sqrt{\Dm}\), \(g_+=\sqrt{\Ag}\).
The last root may be zero.

For \(a_*=|a|\ge(x-1)m\),
\(\sqrt\Delta>a_*-3>xm/2\). Hence
\[
f_+>\frac{xm}{2}C_0^2,\qquad f_+>3C_0^2,\qquad
K(xm^2-1)<\tau_+<K(xm^2+1).
\]
The numerical lower bounds imply
\begin{equation}
\tau_+>2C_0,\quad f_+>2\tau_+,\quad
Q>\Df>\Dt>1,\quad Q>C_0^2.
\label{eq:gaps}
\end{equation}
For instance \(f_+>2\tau_+\) follows from
\(mK(m-1)^2>4(m^2+1/x)\).
For \(b=cT-R\), two cases give
\begin{equation}
\begin{array}{ll}
b=0:&C_0<y_+<C_0+1,\\
b\ne0:&|V|\ge f_+^2-C_0>2f_+,\quad y_+>|V|>2f_+.
\end{array}
\label{eq:ycase}
\end{equation}
In the first case the correction
\((C_0^2-1)/(Q-1)\) is less than \(1\).

The main root satisfies
\begin{equation}
\frac{C_0}{5}<\frac{d_+}{|H|}<\frac{C_0}{2}.
\label{eq:mainbound}
\end{equation}
Use \(\sqrt\Delta>a_*-3\), \(|H|\le4a_*+3\) for the lower
bound, and \(d_+<(a_*+3)C_0\), \(|H|\ge4a_*-3\) for the upper.
The determinant identity
\begin{equation}
12G=4q^2-\bigl(3(U+Z_c)-4q\bigr)^2-3(U-Z_c)^2
\label{eq:detbound}
\end{equation}
gives \(0<G\le q^2/3<q^2\). Since
\(\eta\zeta\le k^2/4\), the positive subtracted term in \eqref{eq:guard}
implies
\begin{equation}
g_+<\frac{|q|K}{2}<\frac{C_0}{10}.
\label{eq:guardbound}
\end{equation}
The last step uses \(m\ge|q|^3\) and \eqref{eq:cstrip}.
The huge determinant coefficients disappear from this estimate.

\begin{proposition}[Strict separation]\label{prop:separation}
After omitting a zero \(g_+\), each entry in either order below exceeds
the sum of all earlier entries:
\[
\begin{array}{ll}
b=0:&g_+,\ d_+/|H|,\ y_+,\ \tau_+,\ f_+,\\
b\ne0:&g_+,\ d_+/|H|,\ \tau_+,\ f_+,\ y_+.
\end{array}
\]
\end{proposition}
\begin{proof}
The two small roots have sum less than \(0.6C_0\), with
\(d_+/|H|>2g_+\). For \(b=0\),
\[
y_++d_+/|H|+g_+<1.6C_0+1<2C_0<\tau_+,
\]
and \(y_+>C_0\), \(f_+>2\tau_+\).
For \(b\ne0\), successively use
\(\tau_+>2C_0\), \(f_+>2\tau_+\), \(y_+>2f_+\).
\end{proof}

\section{Recovering all integer witnesses}\label{sec:kernel}

\begin{lemma}[Square classes]\label{lem:radicals}
If positive numbers with rational squares have strictly superincreasing
magnitudes, and a signed sum is rational, then each is rational.
\end{lemma}
\begin{proof}
In the multiquadratic extension they generate, a nonrational root is
changed in sign by some Galois automorphism. Subtract its image of
the rational sum from that sum. The resulting nonempty signed
subsum cannot vanish: its largest magnitude exceeds the sum of all
others. Independence of the original square classes is not assumed.
\end{proof}

\begin{lemma}[Adjacent-coefficient descent]\label{lem:descent}
If \(Q>1,V\) are integers and \(y\in\Q\), then
\[
QV^2-(Q-1)y^2=1\quad\Longrightarrow\quad y\in\Z.
\]
\end{lemma}
\begin{proof}
Take absolute values and write \(y=a/d\) in lowest terms.
The equation implies \(d^2\mid Q-1\). At \(|V|=1\), \(|y|=1\).
For \(V,y>0\), \(V>1\), use
\[
V'=(2Q-1)V-2(Q-1)y,\qquad y'=(2Q-1)y-2QV.
\]
The norm is unchanged, \(V'\) is integral, and the reduced
denominator of \(y'\) remains \(d\), since \(2Q-1\equiv1\bmod d\).
The inequalities \(y>V\) and \((Q-1)y<QV\) give
\(-V<V'<V\). The norm excludes \(V'=0\).
Descent reaches \(|V|=1\) without changing the denominator.
\end{proof}

\begin{definition}[Normalized integer system]\label{def:kernel}
Let \(\Cal\) use the same registers, omit \(n\), and add the five
integer coordinates \(f,\tau,y,\gamma,p_1\), with equations
\begin{align}
f^2&=\Df,&\tau^2&=\Dt,\nonumber\\
QV^2-(Q-1)y^2&=1,&
(b_0+H\gamma)^2&=\Dm,&\Ag&=p_1^2.
\label{eq:kernel}
\end{align}
This is a seventeen-integer system, with no rational notation in
its defining equations.
\end{definition}

\begin{theorem}[New algebraic reduction]\label{thm:algebra}
\(P_{13}\) has an integer zero if and only if \(\Cal\) has an
integer solution. This assertion and proof are formalizable in ZFC.
\end{theorem}
\begin{proof}
At a zero, Proposition~\ref{prop:guard} gives the guards, then
Proposition~\ref{prop:separation} gives separation.
Divide the vanishing factor by \(H\). Lemma~\ref{lem:radicals}
makes its five roots rational.
The signed roots \(d,f,\tau,p_1\) of the integer radicands are
integers; Lemma~\ref{lem:descent} makes the signed auxiliary root
\(y\) integral. Absorb the factor's signs into these roots:
\[
H(n+V)+b_0=d+H(f+\tau+p_1+y).
\]
Therefore
\begin{equation}
\gamma=n+V-f-\tau-p_1-y=(d-b_0)/H\in\Z.
\label{eq:gamma}
\end{equation}
This restores \eqref{eq:kernel}.
Conversely, given its witnesses, set
\begin{equation}
n=\gamma+f+\tau+p_1+y-V.
\label{eq:choose-n}
\end{equation}
The guard gives \(c\ne0,\delta>0\), and the specified signed roots
make one factor zero. Every quantity is chosen before \(n\);
there is no circular witness selection.
\end{proof}

\subsection{Restoration of the original strong kernel}

Let \(i=2r+1\), and define
\begin{align}
I_{\rm old}&=\Delta i c^2,&r_{\rm old}&=(\Delta i-1)/2,\nonumber\\
o_{\rm old}&=f(cT-R),&
j_{\rm old}&=\Df T-1-R\Delta i^2c^3.
\label{eq:restore}
\end{align}
The first quotient is integral because \(\Delta i\) is odd.
Using \(f^2=\Df\), polynomial identities give
\begin{align}
I_{\rm old}&=(2r_{\rm old}+1)c^2,&
I_{\rm old}^2&=\Delta(f^2-1)=Q,\nonumber\\
fo_{\rm old}-c&=V,&cj_{\rm old}-R&=V.
\label{eq:restore2}
\end{align}
This restores both old quotients without an inferred divisibility
shortcut. The first centered parent root is \(\tau-XY^2k\);
the main root is \(b_0+H\gamma\).
The transport row vanishes identically:
\[
(K_0+W)(q-F-\alpha)-F-z(q-1)=0.
\]
The old radix coordinate is \(J_3=2j\), and \(\Ag=p_1^2\)
restores exactly the joint guard with \(\pos(p_1,p_2,p_3)\).
These are the incoming signed, residue-oriented,
ordinary-strong equations~\cite{joint,residue}.
The new normalization chooses a subfamily of strong witnesses.
Completeness uses the canonical regeneration proved next, rather
than claiming that every old witness is already normalized.

\section{Canonical completeness and consistency calibration}\label{sec:complete}

The preceding restoration proves soundness into the inherited signed
kernel. Completeness requires a new canonical choice of the strong and
auxiliary witnesses. It does not require every old solution to admit
this normalization with its original auxiliaries.

Write the ordinary Pell sequences as
\[
\chi_A(s)+\psi_A(s)\sqrt{A^2-1}
 =(A+\sqrt{A^2-1})^s,\qquad s\in\N.
\]
The inherited canonical native witness has \(A=a+2>1\) even,
\(R\equiv3\pmod4\), and
\[
c=\psi_A(R)>0,\qquad D=\chi_A(R).
\]
The coefficient proof gives \(q=2^e\) with \(e\) odd, odd
\(Y/q^3\), and the residue-oriented scale \(X/(4q)\equiv1\pmod3\).
Also \(J=(q-1)/\beta\) is odd and \(1\bmod3\), so \(J=6j+1\).
Thus the stronger radix parity in \eqref{eq:front1} preserves canonical
existence~\cite{coeff,residue}.

\begin{lemma}[Odd normalized strong witness]\label{lem:normal}
Put
\[
I_*=\psi_A(Rc),\qquad f=\chi_A(Rc),\qquad i=I_*/c^2.
\]
Then \(i\) is a positive odd integer and
\(f^2=1+\Delta i^2c^4\).
\end{lemma}
\begin{proof}
The composition identity gives \(I_*=c\psi_D(c)\).
Expanding \((D+c\sqrt\Delta)^c\), the coefficient of
\(\sqrt\Delta\) is
\[
I_*=\sum_{\substack{1\le s\le c\\s\text{ odd}}}
 \binom cs D^{c-s}c^s\Delta^{(s-1)/2}.
\]
The \(s=1\) term is \(c^2D^{c-1}\), and every other term
has a factor \(c^3\). Hence \(c^2\mid I_*\).
Since \(A\) is even, the recurrence
\(\psi_A(s+2)=2A\psi_A(s+1)-\psi_A(s)\) gives
\(\psi_A(s)\equiv s\pmod2\).
Consequently \(c\) and \(I_*\) are odd, and division by the odd
integer \(c^2\) makes \(i\) odd. This parity conclusion does not
rely on the parity of the leading binomial term.
The Pell identity gives the displayed norm.
\end{proof}

Set \(r=(i-1)/2\) and \(S=\Delta I_*\).
The inherited squared-modulus auxiliary lemma~\cite{squared} gives
\[
V_*=\frac{\chi_S(R)}S\in\Z_{>0},\qquad y=\psi_S(R)>0,
\]
with
\[
S^2V_*^2-(S^2-1)y^2=1,\qquad
V_*\equiv-c\pmod{f^2},\qquad V_*\equiv-R\pmod c.
\]
Its hypotheses are \(R\equiv3\pmod4\), \(c=\psi_A(R)\),
\(f^2-\Delta I_*^2=1\), and \(c\mid I_*\), all established above.
Since \(f^2\equiv1\pmod c\), \(\gcd(c,f)=1\), and therefore
\[
T=\frac{V_*+c+Rf^2}{cf^2}\in\Z.
\]
It gives \(V_*=f^2(cT-R)-c\), exactly \eqref{eq:front3}.
The main native coordinates and their determinant/strip guards remain
unchanged. Choose their joint positivity representation and then choose
\(n\) by \eqref{eq:choose-n}. This completes the canonical reverse
construction.

\begin{theorem}[Ordinary consistency equivalence]\label{thm:consistency}
Use the fixed tuple in Appendix~\ref{app:coeff} and the pinned ordinary
native, signed-domain, and K215 calibration theorems. Then
\[
\mathrm{ZFC}\ \vdash\
\left[
  \bigl(\forall \mathbf x\in\Z^{13}\;
             P_{13}(\mathbf x)\ne0\bigr)
  \ \longleftrightarrow\ \Con(\mathrm{ZFC})
\right].
\]
The same assertion holds for each of the companion polynomials
below, with its stated integer dimension.
\end{theorem}
\begin{proof}
The new argument restores an inherited source witness from every zero.
The canonical construction above produces a zero from every canonical
native witness. The ordinary source theorem identifies existence of
native witnesses with blank halting of the fixed K215 table, and the
ordinary K215 theorem identifies that halting statement with existence
of a finite syntactic ZFC refutation~\cite{coeff,calibration}.
These arguments use ordinary finite coding, induction, Pell arithmetic,
first-order completeness, and the inherited set-theoretic calibration,
all as mathematics in ZFC. Compose the equivalences and negate them.
\end{proof}

The dependencies in this theorem are explicit imported mathematical
results. In particular, no inference of internal provability is made
merely from the external Lean statement
\[
\texttt{BlankHalts machine215}\ \longleftrightarrow\
 \texttt{BProv ZFCax\_s [] Form.fBot}.
\]
The ordinary K215 exposition itself states a ZFC-provable equivalence.
Its external Lean proof passes through the numerical calculus, ZF
models, and constructible-choice calibration. The repository's
machine-checked translation to its actual sealed
\texttt{Form} consistency sentence remains open~\cite{boundary}.
The archive's \texttt{provenance/CALIBRATION\_STATUS.md} separates
these claims. If the imported ordinary calibration is treated as an
obligation to be re-established, the new unconditional contribution is
the exact source-zero equivalence; the final consistency interpretation
then depends on that calibration. Neither bounded testing nor
proof-assistant soundness is substituted for the missing proof artifact.

\section{Exact degrees of the five-root equations}\label{sec:degree}

Every fixed numeral has degree zero, and each supplied unknown has
degree one. In the base coordinates,
\[
\begin{array}{c|rrrrrrrrr}
\text{register}&H&\Delta&R&\Df&Q&V&C_y&\Dt&\Ag\\\hline
\text{degree}&6&12&5&18&30&23&76&34&44 .
\end{array}
\]
These are degrees after the literal substitutions, not degrees on the
solution set.

\subsection{The unshifted degree 1408}
The four integer radicands in \eqref{eq:main} have degrees
\(14,30,46,56\); the auxiliary numerator \(H^2C_y\) has degree
\(88\), while \(\deg\delta=30\). The universal norm is weighted
homogeneous of weight \(32\), with weight one for the center and
weight two for each radicand. In the unshifted equation the center
has degree seven. A term using the auxiliary radicand to power \(s\)
has degree at most
\[
30(16-s)+88s+56(16-s)=1376+2s\le1408.
\]
Only \(s=16\) reaches \(1408\), and its coefficient is \(+1\).

To show that this term survives actual coefficient specialization, put
\begin{align}
q_0&=6\beta j,&W_0&=12w,&Y_0&=2vq_0^3,&X_0&=q_0W_0,\nonumber\\
a_0&=X_0Y_0,&i_0&=2r,&F_2&=W_0C-q_0z,&R_0&=-q_0^3F_2.
\label{eq:topforms}
\end{align}
Then
\[
(H^2C_y)_{88}=16a_0^{10}i_0^6c^{12}R_0^2,
\]
a nonzero polynomial for the authentic fixed \(\beta>0\).
The unshifted degree-\(1408\) part is its sixteenth power.

\subsection{The translated degree 1392}
For the circuit of Section~\ref{sec:circuit}, use
\[
D=\Dm,\quad B_g=H^2\Ag,\quad
A=H^2\Df,\quad B=H^2\Dt,\quad \Cn=H^2C_y.
\]
Give the two formal root symbols \(u,v\) weights \(7,28\).
A coefficient on \(1,u,v,uv\) then has its ordinary degree plus
\(0,7,28,35\). Multiplication respects this filtration.

For \(s=Hn+b_0\), the center is \(x=s+HV\).
The scalar part of \(J_N\), evaluated at \(x+u+v\), has the exact
cancellation
\[
(J_N)_0=\delta(s^2+2HVs+D+B_g-A-B)-H^2(V^2-1).
\]
Its degree is \(86\), led by \(\delta B_g\).
The four weighted component degrees are
\[
86,\quad66,\quad87,\quad65.
\]
The unique largest component is \(2Q_0x_0v\), of weight \(87\).
Its square has leading scalar \((2Q_0x_0)^2(B_g)_0\), of degree
\(174\); the correction in \(K_N\) has degree at most \(164\).
The correction in \(N_3\) has weight at most \(312\), so \(N_3\)
has unique leading scalar
\[
16Q_0^4x_0^4(B_g)_0^2
\]
of degree \(348\). Each of its four conjugates has this same
leading scalar. Their product therefore has exact degree
\(4\cdot348=1392\).

Explicitly,
\[
\eta_0=-hX_0Y_0^2,\quad
Q_0=a_0^4i_0^2c^4,\quad
V_0=-a_0^2i_0^2c^4R_0,\quad
x_0=4a_0V_0,\quad (B_g)_0=-16a_0^2\eta_0^4.
\]
Thus
\[
(P_{13})_{1392}
=2^{80}a_0^{128}i_0^{64}c^{128}R_0^{16}\eta_0^{32}\ne0.
\]
For a directly checkable specialized certificate,
\[
\left[j^{800}w^{160}v^{192}h^{32}r^{64}c^{128}z^{16}\right]P_{13}
 =2^{1456}3^{960}\beta^{800}>0.
\]
This proves exact degree, rather than relying on cancellation-free
gate estimates or modular interpolation.

\input{report/crt1088.tex}

\section{Four-root and three-root companions}\label{sec:companions}

Define \(\F_{16}\) and \(\F_8\) by the corresponding products
over four and three signed radicals. Retain the base registers
\eqref{eq:front1}--\eqref{eq:guard}, and write
\[
J_g=\Ag-p_1^2,\qquad
\ell=XY^2k,\qquad F_{\rm s}=(a+2)I.
\]
A useful finalizer is
\[
\mathcal P=N+S+NS,\qquad S\text{ a sum of integer squares}.
\]
For integer inputs, \(\mathcal P=0\) is equivalent to \(S=N=0\):
\((1+S)(1+N)=1\) and \(1+S\) is a positive integer.

\subsection{Fourteen integers, degree 696}
Add \(p_1\) to the thirteen base coordinates and set
\begin{align}
Z_{14}&=H(n+2V+2\ell+k+2F_{\rm s})+2b_0,\nonumber\\
N_{14}&=\delta^8\F_{16}\!
 \left(Z_{14};4\Dm,4H^2\Df,4H^2\Dt,\frac{4H^2C_y}{\delta}\right),
 \nonumber\\
S_{14}&=J_g^2,\qquad
P_{14}=N_{14}+S_{14}+N_{14}S_{14}.
\label{eq:four}
\end{align}
The literal polynomial has \(150\) operations.
At a zero, the guard and four-root separation give integral
\(d,f,\tau,y\). A signed factor yields
\(Z_{14}=2d+2H(f+\tau+y)\).
Hence \(H\mid2(d-b_0)\), and oddness of \(H\) gives
\(\gamma=(d-b_0)/H\in\Z\).
Conversely, choose
\[
n=2\gamma+2f+2\tau+2y-2V-2\ell-k-2F_{\rm s}.
\]
This is an existence-preserving half-centering, with no parity
assumption on the independent coordinate \(n\).

Here is a precise degree argument. Replace all independent variables
by a common indeterminate multiple and work at infinity in the
formal Laurent-series field. The three centered radicals have
centers and corrections
\[
\begin{array}{c|c|c}
\text{center}&\text{radicand minus center squared}&
 \text{degrees: center, correction}\\\hline
2HF_{\rm s}&4H^2(1-I^2)&15,18\\
H(2\ell+k)&H^2(4-k^2)&23,26\\
2HV&4H^2(V^2-1)/\delta&29,28 .
\end{array}
\]
Their root deviations have degrees at most \(3,3,-1\).
Of the sixteen conjugate factors, eight first fail to cancel the
last center, four first fail to cancel the middle center, and two
first fail to cancel the first center. Their degrees are
\(29,23,15\), respectively. The final main-root pair has degree
\(14\), with leading form \(H_0^2n(n+c)\), where \(H_0=4a_0\).
Clearing \(\delta^8\) adds \(240\). Therefore
\[
\deg N_{14}=8(29)+4(23)+2(15)+14+240=608.
\]
Its highest form is
\[
(N_{14})_{608}
 =2^{60}a_0^{66}i_0^{34}c^{68}R_0^8\ell_0^4n(n+c),
 \qquad \ell_0=hX_0^2Y_0^3.
\]
The guard square has degree \(88\), led by \(\eta_0^8\).
Consequently \(\deg P_{14}=696\), certified, for example, by
\[
[j^{396}w^{82}v^{94}h^{12}r^{34}c^{68}z^8n^2]P_{14}
 =2^{748}3^{478}\beta^{396}>0.
\]

The unshifted fourteen-variable companion replaces the center and
radicands by
\[
Hn+b_0;\quad \Dm,H^2\Df,H^2\Dt,H^2C_y/\delta .
\]
Its cleared norm has degree \(704\), led by \((H^2C_y)_{88}^8\);
the same finalizer has exact degree \(792\) and \(137\) operations.
The four-root argument recovers the main residue directly.

\subsection{Fifteen integers, degree 344}
Add \(\gamma,p_1\) to the thirteen base coordinates and use
\begin{align}
Z_{15}&=n+2V+2\ell+k+2F_{\rm s},\nonumber\\
N_{15}&=\delta^4\F_8(Z_{15};4\Dt,4\Df,4C_y/\delta),\nonumber\\
S_{15}&=((b_0+H\gamma)^2-\Dm)^2+J_g^2,\nonumber\\
P_{15}&=N_{15}+S_{15}+N_{15}S_{15}.
\label{eq:fifteen}
\end{align}
There are \(131\) operations.
The two explicit residuals provide the main and guard rows.
Three-root separation supplies integral \(f,\tau,y\), and conversely
\[
n=2(f+\tau+y)-2V-2\ell-k-2F_{\rm s}
\]
makes a norm factor vanish.

The Laurent centers \(2V,2\ell+k,2F_{\rm s}\) have degrees
\(23,17,9\). Their deviations have degrees at most
\(-7,-3,-3\), so the final fully canceled factor has leading term
\(n\), degree one. The eight-factor norm therefore has degree
\[
4(23)+2(17)+9+1+4(30)=256.
\]
Its leading form is
\(2^{14}Q_0^4V_0^4\ell_0^2(F_{\rm s})_0n\ne0\).
The guard square again has degree \(88\), giving exact final
degree \(256+88=344\).
These degree computations are formal polynomial calculations;
they do not use numerical asymptotics of integer solutions.

\section{Four quotient-fusion companions}\label{sec:quotient}

The archive also gives less radical-intensive equations with
variable/degree pairs
\[
(20,24),\quad(19,28),\quad(18,32),\quad(17,88).
\]
They use the same six coefficients. Their full literal residual
lists and independent certificates are in
\texttt{companions/quotient\_fusion/PROOF.md} and its JSON files.

For completeness, the twenty-variable presentation is specified here.
Its independent coordinates are
\[
j_3,F,\alpha,z,f,h,r,v,w_3,\tau,\eta,y,\gamma,T,a,c,t,p_1,p_2,p_3.
\]
Use \(J=3j_3+1\), \(q=\beta J+1\), \(W=4(3w_3+1)\),
\(X=qW\), \(Y=(2v+1)q^3\), \(C=q-F-\alpha\), and
\begin{align*}
E&=a-Y,& H&=4a+3,&\Delta&=a^2+H,\\
R&=[q(q-F)-C](q^2-1)+(M_C+qM_F)J,&&&
k&=R+1+hE.
\end{align*}
Set \(I=(2r+1)c^2\), \(Q=\Delta(f^2-1)\),
\(B=f^2(cT-R)\), \(\zeta=k-\eta\), and retain the determinant
\(G\) of \eqref{eq:guard}. Sum the squares of these eight residuals:
\begin{align*}
&a-(X+1)Y, &&c-\eta-Yk,\\
&\tau^2+(2\tau-k)E(c-\eta)-1,
&&(X+ac+\gamma H)^2-\Delta c^2-1,\\
&(K_0+W)C-F-z(q-1),&&I^2-Q,\\
&(B+y)^2-Q(c+y)(2B+y-c)-1,\\
&G\eta\zeta-(t^2-G-\eta\zeta+1)^2-\pos(p_1,p_2,p_3).
\end{align*}
This polynomial has degree \(24\) and \(120\) operations.
It centers the old first and auxiliary roots by
\[
\tau_{\rm old}=E(c-\eta)+\tau,\qquad y_{\rm old}=B+y.
\]
The auxiliary identity is
\(y_{\rm old}^2-(B-c)^2=(c+y)(2B+y-c)\).

Replacing \(F,\alpha,z\) by \(C,z\) through the monic transport
parameterization gives nineteen variables, degree \(28\), and
\(116\) operations. Next substitute \(a=(X+1)Y\), \(E=XY\),
remove the scale row, and use \(I^2\) instead of \(Q\) in the
auxiliary row. The strong residual makes this row replacement
reversible on zeros. This gives eighteen variables, degree \(32\),
and \(112\) operations. Finally substitute \(\eta=c-Yk\) and
delete its row, obtaining seventeen variables, degree \(88\),
and \(110\) operations.

Soundness of quotient fusion must derive \(c\mid f^2-1\) before
restoring the removed quotient. The inherited main classification
and ordinary strong-rank theorem give
\(c=\pm\psi_A(p)\), \(|f|=\chi_A(m)\), \(p\mid m\), hence
\(c^2\mid f^2-1\). Only then define
\[
o=f(cT-R),\qquad
j_{\rm old}=f^2T-1-R(f^2-1)/c.
\]
These are integers with \(fo-c=B-c=cj_{\rm old}-R\).
Completeness uses the squared-modulus canonical auxiliary choice
already described, with the ordinary \(\Delta\)-scaled odd strong
multiplier. This avoids the invalid shortcut of canceling
\(\Delta\) from the ordinary strong equation.

The exact degree witnesses include
\[
[f^8F^4j_3^{12}]P_{20}=(3\beta)^{12},\qquad
[f^8z^4j_3^{16}]P_{19}=(3\beta)^{16},
\]
and the highest guard forms certify the eighteen- and
seventeen-variable endpoints. The independent sparse verifier
retains symbolic fixed ports until specializing them to the
authentic tuple.

\section{Reproduction and evidence}\label{sec:reproduce}

From the archive root, with Python 3.10 or later, run:
\begin{lstlisting}
python3 build_equations.py
python3 verification/verify_release.py
python3 -O verification/verify_release.py
pdflatex -interaction=nonstopmode -halt-on-error \
  -output-directory=report report/report.tex
pdflatex -interaction=nonstopmode -halt-on-error \
  -output-directory=report report/report.tex
\end{lstlisting}
The release verification uses explicit checks that remain active under
\texttt{-O}. Its results bind the complete numeral payloads,
literal coefficient recipes, declared variables, live arithmetic gates,
exact norm identity, degree certificates, and companion equations.
The editable \TeX\ files and the PDF are both delivered.

The finite evidence and the unbounded proof have different roles.
Exact coefficient dictionaries establish universal polynomial identities.
A nonzero symbolic highest form establishes degree for the actual
fixed slice. Hashes and reconstruction bind the coefficient data.
Sample evaluations can detect transcription mistakes, but cannot prove
that all zeros decode, that a compiler is complete, or that the
consistency biconditional holds. Those assertions use the mathematical
arguments and imported ordinary calibration identified above.

The independent audits were independently implemented or derived within
the research collaboration. They are not a claim of external human
peer review. No full expansion of the gigantic numerical polynomial
is needed: its literal division-free circuit and fixed integer leaves
specify the ordinary polynomial exactly.

\section{What is new and what remains to improve}\label{sec:novelty}

The principal new composition has five parts: monic transport
coordinates, an odd normalized strong witness, squared-modulus
auxiliary restoration, recovery of the main Pell residue through the
norm center, and absorption of the guard through unconditional
realness followed by automatic magnitude separation.
The last step removes the guard's square witness without the enormous
separating weights required by a generic square-condition combiner.
The CRT companion further recovers the odd cubic scale and the main
residue from one integral carrier; its weaker scale is safe because
the scaled guard excludes the new zero-scale boundary first.

Several ingredients are inherited: the radical norm, joint-positivity
lemma, determinant cone, adjacent-coefficient descent, Pell
congruences, and native/refutation calibration. We claim the displayed
new reductions and literal measured endpoints, not worldwide priority
for their general ingredients. The proposal remains K215-based through
its fixed coefficient tuple. The request did not require avoiding
that source, and the structural simplifications do not depend on
Jones's universal polynomial.

There is no established global minimum under all reasonable size
measures. Fewer integer variables can demand much higher degree;
smaller circuits can have much larger coefficients. The degree-\(1088\)
form and the \(186\)-operation degree-\(1392\) form express that
tradeoff concretely. A useful next target is a square-extraction
argument that merges another native coordinate without losing the
guard's quantitative information. Any such attempt must preserve
canonical completeness, not merely the forward decoding direction.
A separate substantial task is to compose the literal native instance
and this reduction in the repository's sealed object-ZFC calculus.

\appendix
\section{Complete bounded coefficient recipe}\label{app:coeff}

The following is the full source-34 finite arithmetic recipe~\cite{coeff}.
All ranges are bounded constants. The ordered \(943\) tile records
\((x_{i0},x_{i1},x_{i2},x_{i3})\), in west, north, east, south order,
are supplied in
\texttt{coefficients/recipe/wang943\_tiles.json}.
Their ordering is part of the definition: selector \(0\) is Start
and selector \(1\) is Forbidden.
The supplied \texttt{k215.tm} contains all \(430\) working actions
from which the source builder reconstructs the whole ordered list.
Its SHA-256 is
\begin{center}\small\texttt{
fdbcef0825886035c3091c8537a69535bc96f9e265f0ae45863543d40ae70090}.
\end{center}
Thus this finite table is delivered data, not an unspecified machine
or a free coefficient parameter.

Put
\[
b=121,\quad L_0=8801,\quad d=bL_0=1064921,\quad
\rho=2^{121},\quad B_0=2^{1064921}.
\]
For \(0\le s<4\), \(0\le j<10\), define
\[
e_s(j)=o_s+2j,\qquad(o_0,o_1,o_2,o_3)=(1020,1000,980,960),
\]
and \(w_j=2^{10+2j}\) for \(0\le j<44\). Let
\[
\begin{split}
\mathcal E={}&\{0,\ldots,942\}
 \cup\{e_s(j):0\le s<4,\ 0\le j<10\}\\
 &\cup\{1063,1064,1095,1096,2195,3296\}.
\end{split}
\]
It has \(989\) elements and maximum \(3296\).
For each tile \(0\le i<943\), set
\[
c_i=1+\sum_{s=0}^3\sum_{j=0}^9
       \operatorname{bit}_j(x_{is}+512)w_{10s+j}
       +\sum_{j=40}^{43}w_j .
\]
Here \(\operatorname{bit}_j(m)=\lfloor m/2^j\rfloor\bmod2\)
is used only to define the finitely many fixed integer coefficients.
It is not an operation on a quantified Diophantine variable.
The other coefficients are
\[
c_{e_s(j)}=w_{10s+j},\quad
(c_{1095},c_{1096},c_{2195},c_{3296})=(w_{40},w_{41},w_{42},w_{43}),
\quad c_{1063}=c_{1064}=0.
\]
Define
\begin{align*}
\mu&=1022+\sum_{j=0}^{43}w_j,\\
Z_{\rm pad}&=(2^{119}-1)\sum_{e=8792}^{8798}\rho^e
                +(2^{27}-1)\rho^{8799},\\
D_C&=\rho^{40}+\rho^{41}+\rho^{1100}+\rho^{2200}
       +\sum_{e\in\mathcal E}c_e(\rho^{4393-e}+\rho^{5495-e}),\\
D_R&=\rho,\\
M_C&=B_0-1-\sum_{\substack{e\in\mathcal E\\e\ne1}}\rho^e
                   -2\rho^{1063},\\
M_{F,\rm native}
 &=4+\mu(\rho^{4393}+\rho^{5495})
     +\sum_{j=0}^9(\rho^{1000+2j}+\rho^{1021+2j})
     +\rho^{2195}+\rho^{3296}+Z_{\rm pad}.
\end{align*}
The six actual fixed ports are precisely
\[
\boxed{\begin{gathered}
\beta=B_0-1,\quad K_0=D_C+B_0D_R,\quad M_C,\quad
M_F=M_{F,\rm native}+B_0-1,\\
A_{\rm cone}=2^{1113},\qquad B_{\rm cone}=2^{666104}.
\end{gathered}}
\]
All displayed exponents are fixed integers. In particular,
\(M_F\) already includes its shift; shifting it again would change
the equation and invalidate the binding.

\begin{table}[ht]\centering
\begin{tabular}{lrr}\toprule
Fixed numeral & Bit length & Minimal unsigned bytes\\\midrule
\(\beta\)&1\,064\,921&133\,116\\
\(K_0\)&1\,065\,043&133\,131\\
\(M_C\)&1\,064\,921&133\,116\\
\(M_F\)&1\,064\,922&133\,116\\
\(A_{\rm cone}\)&1\,114&140\\
\(B_{\rm cone}\)&666\,105&83\,264\\\midrule
Total&&615\,883\\\bottomrule
\end{tabular}
\caption{Exact fixed payload sizes before compression.
The archive stores deterministic gzip encodings of the minimal
unsigned little-endian bytes and binds their raw hashes.}
\end{table}

The source proof verifies the whole coefficient interface, not just
these sizes. It establishes local tile decoding, carry-free geometry,
whole-cell rotation, mask populations, positive canonical slack, and
an explicit rank controller. In particular
\[
\operatorname{pc}(M_C)+\operatorname{pc}(M_{F,\rm native})=d,\quad
M_C\equiv2\pmod4,\quad M_{F,\rm native}\equiv4\pmod8.
\]
The controller uses
\(d=11^2\cdot13\cdot677\), \(\lambda=2341\), and a proved Boolean
subset coverage lemma; it does not assume that \(d\) is a power of
five. Its construction supplies odd canonical exponents, odd cubic
scale, and the parity promises used in Lemma~\ref{lem:normal}.
The complete inherited proof and literal builder are included,
with their internal source pins and licenses preserved.

\section{Archive map and integration}\label{app:files}

\begin{longtable}{@{}>{\raggedright\arraybackslash}p{0.38\linewidth}>{\raggedright\arraybackslash}p{0.58\linewidth}@{}}
\toprule
Location & Contents\\\midrule
\endhead
\texttt{report/}&
This editable report, the CRT section, and the compiled PDF.\\
\texttt{build\_equations.py},
\texttt{five\_square\_norm.py}&
Literal equation generator and division-free universal norm evaluator.\\
\texttt{equations/}&
The six norm circuits, declared variable lists, exact degree metadata,
and operation ledgers.\\
\texttt{coefficients/}&
All six fixed numeral payloads, binding manifest, complete bounded
recipe, machine table, and ordered tile list.\\
\texttt{companions/quotient\_fusion/}&
The four quotient-fusion equations, source proof, generator, and
independent verifier.\\
\texttt{verification/}&
Integrated release audit, exact norm identity and degree checks,
and machine-readable receipts.\\
\texttt{provenance/}&
Inherited interfaces, calibration status, independent mathematical
audit, and reviewed frontier.\\
\texttt{provenance/upstream\_ordinary/}&
Verbatim ordinary calibration proofs and exact external endpoint
sources, each bound by commit, path, SHA-256, and byte length.\\
\texttt{licenses/}&
Preserved repository, source, and machine notices.\\
\bottomrule
\end{longtable}

The folder is intended for a new proposal directory in the ProveIt
Diophantine-consistency collection. No existing proposal number is
reserved. All runtime paths are relative to the delivered folder.
Keep the whole six-numeral binding together when copying or adapting
a circuit. The coefficients, source system, and formula are a single
fixed instance.

\begin{thebibliography}{12}\small\raggedright
\bibitem{repo}
V. Reshetnikov, \emph{ProveIt}, repository snapshot
\texttt{9a33342a03d19cf5b061b2083c6b5d3f27f752ed},
24 October 2026.
\url{https://github.com/VladimirReshetnikov/ProveIt}.

\bibitem{frontier}
ProveIt, \emph{Con-ZFC Diophantine forms}, research catalogue and
frontier at the stated snapshot. The exact catalogue README is
retained as \texttt{provenance/CON\_ZFC\_FORMS\_README.md}.
\url{https://github.com/VladimirReshetnikov/ProveIt/tree/9a33342a03d19cf5b061b2083c6b5d3f27f752ed/Computability/DiophantineConsistency/con-zfc-diophantine-forms}.

\bibitem{incoming}
ProveIt, \emph{Incoming research packets}, including
\emph{Structural Reductions, Round 11} and
\emph{ZFC Diophantine Joint Guards}, 24 October 2026.
\url{https://github.com/VladimirReshetnikov/ProveIt/tree/9a33342a03d19cf5b061b2083c6b5d3f27f752ed/docs/incoming}.

\bibitem{boundary}
ProveIt, \emph{The object-language ZFC obligation},
\path{Computability/DiophantineConsistency/OBJECT_ZFC.md}.
Verbatim boundary document retained as
\texttt{provenance/FORMALIZATION\_BOUNDARY.md} in this archive.

\bibitem{calibration}
ProveIt, \emph{Ordinary proof of the combined215 machine};
\emph{Ordinary strong norms and signed witnesses for the blank
computation}; and \emph{Native arithmetic and the remaining
calibration contract}. Exact copies of the ordinary proof texts,
source paths, and hashes are in
\texttt{provenance/upstream\_ordinary/SOURCES.json}.
\url{https://github.com/VladimirReshetnikov/ProveIt/blob/9a33342a03d19cf5b061b2083c6b5d3f27f752ed/Computability/BusyBeaver/Research/ZFC/candidate-k215/proof.md}.

\bibitem{joint}
ProveIt incoming packet,
\emph{ZFC Diophantine Joint Guards}, 24 October 2026,
\texttt{joint/PROOF.md} and \texttt{guards/PROOF.md}.
Relevant ordinary interfaces are retained in
\texttt{provenance/JOINT\_GUARD\_AND\_PROJECTIONS.md} and
\texttt{provenance/RESIDUE\_AND\_DETERMINANT.md}.

\bibitem{residue}
ProveIt, source 70 \emph{Residue orientation} and source 71
\emph{Determinant guards}; signed bootstrap and exact canonical
coefficient composition. Pinned files retained in the archive's
provenance directory.

\bibitem{squared}
ProveIt, source 51, \emph{Squared-modulus auxiliary reconstruction},
\texttt{51-sparse-arith/short/SQUARED\_MODULUS\_PROOF.md};
and source 61, positive auxiliary reconstruction lemma.
The squared-modulus proof is retained as
\texttt{provenance/SQUARED\_MODULUS\_PROOF.md}.

\bibitem{coeff}
ProveIt, source 34,
\emph{A mixed-radix, parity-separated, unrounded coefficient tuple},
\texttt{34-integer14/coefficients/COEFFICIENT\_PROOF.md}.
Complete recipe, materialized coefficients, and verbatim proof in
this archive.

\bibitem{sun2026}
Z.-W. Sun, \emph{Ten unknowns for Hilbert's tenth problem over the
integers}, arXiv:2609.01594v2, revised 2 September 2026.
\url{https://arxiv.org/abs/2609.01594}.

\bibitem{sun2021}
Z.-W. Sun, \emph{Further results on Hilbert's Tenth Problem},
Science China Mathematics \textbf{64} (2021), 281--306.
\url{https://arxiv.org/abs/1704.03504}.

\bibitem{jones}
J. P. Jones, \emph{Universal Diophantine equation},
Journal of Symbolic Logic \textbf{47} (1982), 549--571.
\url{https://doi.org/10.2307/2273588}.
\end{thebibliography}
\end{document}


